\section{The Robustness Audit}

Many of the findings reported earlier were computed assuming a clean three-dimensional space. Now that physical
space has been identified as the cubic cusp, a periodic crystal rather than an idealized continuum lattice, the
pressing question is whether those findings survive. The reassuring fact is that most of them do not live on the
flat slice at all. They live in the bulk coin or on the boundary sphere, where the cusp question simply does not
reach. This section sorts every finding by where it lives, checks the ones that genuinely depend on clean
three-space, and lays out the worst worries worst-first.

\subsection{Three layers, three fates}

The findings split into three layers, and each layer has its own relation to the cusp reality.

\begin{itemize}
  \item \emph{Bulk findings}, which are cusp-independent and unaffected. Spin, triality, the $SO(10)$ gauge
  group, $\sin^2\theta_W = 3/8$, the mass relations, the holographic $1/r^2$ correlator, cosmology, the
  hierarchy, and the toolkit. These live in the four-dimensional bulk coin or are pure group theory, so the cusp
  question does not touch them.
  \item Boundary findings on the $S^3$ sphere, also unaffected. Holography and quantum nonlocality belong here,
  and the cusp does not bear on them.
  \item Flat-layer findings, which need clean three-space and are checked here. The Dirac equation, Lorentz
  invariance, $1/r$ gravity, electromagnetism, solitons, the fermionic-matter result, and matter formation. These
  are the only findings that could be threatened by the cusp reality, and all of them hold.
\end{itemize}

\subsection{The check on the cusp}

To check the flat-layer findings, the relevant observables were measured on the $\{4,3,4\}$ cubic cusp and, for
contrast, on a generic aperiodic slice. The two surfaces give clearly different verdicts.

\begin{table}[h]
\centering
\begin{tabular}{lcc}
\toprule
observable & cubic cusp (physical space) & generic slice (bad extraction) \\
\midrule
spectral dimension & $2.97$ (clean 3D) & $2.16$ (sub-3D) \\
isotropy (light-cone CV) & $0.14$ & $0.32$ (rough) \\
\bottomrule
\end{tabular}
\end{table}

The flat-layer physics holds on the cubic cusp and degrades on the generic slice
\cite{pollard2026vibetest}. The reading is twofold. The earlier results survive provided physical space is the
cusp cubic, which the previous sections established that it is. And they would fail on a disordered slice, which
is therefore not physical space. This does not merely confirm the findings, it also confirms the identity
of space, since only the right surface gives the right physics.

\begin{result}[The flat layer holds on the cusp]
On the $\{4,3,4\}$ cubic cusp the flat-layer observables read spectral dimension $2.97$ and light-cone isotropy
$0.14$, while a generic slice degrades to $2.16$ and $0.32$. The flat-layer findings survive on the cusp and fail
on a disordered slice \cite{pollard2026vibetest}.
\end{result}

\subsection{The problem ledger}

The deepest worries about the theory on the real cusp are listed below, ranked worst-first, each with its
current status. This is the worst-first audit rather than a summary of successes.

\begin{table}[h]
\centering
\begin{tabular}{p{0.42\textwidth}p{0.50\textwidth}}
\toprule
problem & status \\
\midrule
the bulk-cusp interface, do the 4D $D_4$ spinors live on the 3D cusp? & Solved. It is a projection,
$SO(4) \to SO(3)$, $4 \to 2+2$, verified. \\
\addlinespace
cubic anisotropy, the cusp lattice breaks rotations & Resolved. Tiny $(E/E_{\mathrm{Planck}})^2$, IR-isotropic,
falsifiable, and smaller than the naive cubic estimate since the projected coin is richer. \\
\addlinespace
dynamical settling, does matter actually gather into clean cusp regions? & Open. A plausible mechanism, solitons
seeking the regular low-energy parabolic sites, and a defined GPU test, but not yet confirmed. The one
genuinely-open robustness item. \\
\addlinespace
absolute quantitative values, masses, couplings, the Born-rule normalization & The frontier, shared with all
such theories. \\
\addlinespace
paracompactness and nonlocality & Resolved. Paracompactness is a feature, the cusp is 3D space, and the
nonlocality is bulk-provided. \\
\bottomrule
\end{tabular}
\end{table}

One point in the ledger deserves emphasis, because it is subtle. The cubic anisotropy is the cusp's, not the
bulk's. The bulk $24$-cell coin is isotropic to order four, the good $D_4$ and $F_4$ structure. But matter
propagates on the cubic $\{4,3,4\}$ cusp, whose exact symmetry is the finite octahedral group, anisotropic at
order four and not protected by the bulk $F_4$. So the one ultraviolet Lorentz residual lives specifically on the
cusp, and it cannot be inherited away from the bulk's protection.

\subsection{The one modification}

The single change that the cusp reality forces is the cusp's small cubic anisotropy, a tiny
$(E/E_{\mathrm{Planck}})^2$ ultraviolet Lorentz violation that becomes isotropic in the infrared. Everything else
is either unaffected, because it lives in the bulk or on the boundary, or it holds, because the flat-layer
measurement confirms it on the cusp. One artifact note belongs here. The finite-patch gravity exponent reads
about $0.67$, which is a finite-size and screening artifact of the bounded patch rather than a physical reading.
The clean $1/r$ result is the difference-method value, not the patch number, and the two should not be confused.


The program is robust to the cusp reality. Only the flat-layer physics depends on clean three-space, and it holds
on the cusp. The bulk findings and the boundary findings are untouched, because the cusp question never reaches
them. The one genuinely-open item is the dynamical settling of matter into clean cusp regions, which has a
plausible mechanism and a defined test but is not yet confirmed. That single open question is stated plainly
rather than hidden, and the rest of the audit closes cleanly.
